% !PRINTED 9
\weberpage{9}

The subsidiary groups do not have the characteristics of a group. For in
order that two elements from $Q_1$ should, on composition, again yield an
element of $Q_1$, we should have to have, say,
\[
a_1 b a_2 b = a_3 b,
\]
hence $a_1 b a_2 = a_3$, $b = a_1^{-1} a_3 a_2^{-1}$. Thus $b$ would,
contrary to the hypothesis, be contained in $Q$.

The name subsidiary group is therefore to be understood only in an
improper sense.

Two divisors $Q$, $Q'$ of $P$ always have the element 1 in common. But
they may also have still other elements in common, and these common
elements form a group, which we shall denote by $R$. For if the elements
$a$ and $b$ belong both to $Q$ and to $Q'$, then on account of the group
character of $Q$ and $Q'$ the same holds of the compositum $ab$; hence $R$
is a group. This common group we call the \gesperrt{greatest common divisor} of $Q$
and $Q'$, or also, with a geometrical allusion, the \gesperrt{intersection of} $Q$
and $Q'$. It follows likewise that the elements which are common to any
number of divisors of $P$ form a group, which we likewise designate as the
greatest common divisor or the intersection of all these groups.

In every group we can form divisors by the following procedure.

One divisor is always the unit element by itself.

If we denote the repeated composition of an element with itself by powers
(with $a^0$ the unit element), then the sequence of elements
\[
1,\ a,\ a^2,\ a^3 \ldots,
\]
which all belong to the group $P$, cannot contain nothing but distinct
elements. If, then, the first element that recurs for the second time is
\[
a^\mu = a^{\mu+\alpha},
\]
it follows that $a^\alpha = 1$ must hold, and that $\alpha$ is the least
positive number that satisfies this condition. Then, if $m = q\alpha$ is
any multiple of $\alpha$, $a^m = 1$, and conversely: whenever $a^m = 1$,
$m$ must be divisible by $\alpha$; for otherwise one could set
$m = q\alpha + \alpha'$, where $\alpha'$ is positive and less than
$\alpha$, and we should have $a^{\alpha'} = 1$, contrary to the
hypothesis,
%CONT>>
